\documentclass[11pt]{article}

\usepackage[margin=1in]{geometry}
\usepackage{amsmath,amssymb,amsthm}
\usepackage{graphicx}
\usepackage{hyperref}

\newtheorem{theorem}{Theorem}
\newtheorem{lemma}{Lemma}
\newtheorem{corollary}{Corollary}

\newcommand{\E}{E}
\newcommand{\F}{F}
\newcommand{\sol}{\operatorname{sol}}
\newcommand{\eps}{\varepsilon}

\title{Dual cotype forces the disjoint property}
\author{Run \texttt{fa\_banach\_001}, lane 7}
\date{2026-06-20}

\begin{document}
\maketitle

\begin{abstract}
Botelho, Garcia, and Miranda ask in Remark 2.8 of
\cite{botelho-garcia-miranda} whether there can be a Banach lattice \(E\)
such that \(E^*\) has cotype \(p\), \(2\le p<\infty\), while \(E\) fails the
disjoint property of order \(p\).  The answer is negative.  If \(E^*\) has
cotype \(p\), then every bounded disjoint sequence in \(E\) is weakly
\(p\)-summable.  Hence no such counterexample exists.
\end{abstract}

\section{Source question}

The source is G. Botelho, L. A. Garcia, and V. C. C. Miranda,
\emph{Disjoint \(p\)-convergent operators and their adjoints},
arXiv:2309.03775.

The relevant passage is Remark 2.8 on page 5 of the source PDF:

\begin{center}
\includegraphics[width=0.95\linewidth]{figures/open_problem_crop.png}
\end{center}

In paraphrase, the authors ask whether dual cotype \(p\) can hold while the
underlying Banach lattice fails the disjoint property of order \(p\).  Their
Definition 2.5 says that \(E\) has the disjoint property of order \(p\) when
every bounded disjoint \(E\)-valued sequence is weakly \(p\)-summable.

\section{Proof intuition}

The only missing bridge is order continuity.  A Banach lattice with finite
cotype cannot contain an order-bounded disjoint sequence whose norms stay
away from zero: all signed finite sums are dominated in modulus by the same
positive order bound, but the cotype inequality would force the \(\ell_p\)
sum of the norms to grow like \(N^{1/p}\).  Applied to \(E^*\), this gives
order-continuity of the dual norm.

Once \(E^*\) has order-continuous norm, every bounded disjoint sequence in
\(E\) is weakly null.  Such a sequence, together with its weak limit \(0\), is
relatively weakly compact.  The source paper's Lemma 2.4 then applies because
the sequence lies in the solid hull of this weakly compact set.

\section{The result}

\begin{lemma}[Finite cotype forces order-continuity]\label{lem:cotype-oc}
Let \(F\) be a Banach lattice.  If \(F\) has cotype \(p<\infty\), then the
norm of \(F\) is order continuous.
\end{lemma}

\begin{proof}
Let \(C\) be a cotype \(p\) constant for \(F\).  Suppose that the norm of
\(F\) is not order continuous.  By the standard order-continuity criterion for
Banach lattices, there are \(0\le y_n\le y\) with the \(y_n\)'s pairwise
disjoint and, after passing to a subsequence, \(\|y_n\|\ge \delta>0\) for all
\(n\).

For every finite set \(\{1,\dots,N\}\) and every choice of signs
\(\varepsilon_n=\pm1\), disjointness gives
\[
  \left|\sum_{n=1}^N \varepsilon_n y_n\right|
  =\sum_{n=1}^N y_n
  \le y,
\]
and therefore \(\|\sum_{n=1}^N \varepsilon_n y_n\|\le \|y\|\).  Applying the
cotype inequality to \(y_1,\dots,y_N\) yields
\[
  \delta N^{1/p}
  \le \left(\sum_{n=1}^N \|y_n\|^p\right)^{1/p}
  \le C\left(\int_0^1
        \left\|\sum_{n=1}^N r_n(t)y_n\right\|^2\,dt
      \right)^{1/2}
  \le C\|y\|.
\]
Letting \(N\to\infty\) gives a contradiction.
\end{proof}

\begin{theorem}\label{thm:dual-cotype-disjoint-property}
Let \(E\) be a Banach lattice and let \(2\le p<\infty\).  If \(E^*\) has
cotype \(p\), then \(E\) has the disjoint property of order \(p\).
\end{theorem}

\begin{proof}
By Lemma \ref{lem:cotype-oc}, the Banach lattice \(E^*\) has order-continuous
norm.  A standard characterization of order-continuity of the dual norm says
that every norm-bounded disjoint sequence in \(E\) is weakly null; this is the
criterion cited as Theorem 2.4.14 of Meyer-Nieberg in the source paper.

Let \((x_n)\) be any bounded disjoint sequence in \(E\).  The preceding
paragraph gives \(x_n\to0\) weakly.  Hence
\[
  W=\{0\}\cup\{x_n:n\in\mathbb N\}
\]
is weakly compact.  The sequence \((x_n)\) is contained in the solid hull
\(\sol(W)\), and it is disjoint.  Lemma 2.4 of
\cite{botelho-garcia-miranda} states that, when \(E^*\) has cotype \(p\),
every disjoint sequence contained in the solid hull of a relatively weakly
compact subset of \(E\) is weakly \(p\)-summable in \(E\).  Therefore
\((x_n)\) is weakly \(p\)-summable.

Since \((x_n)\) was an arbitrary bounded disjoint sequence, \(E\) has the
disjoint property of order \(p\).
\end{proof}

\begin{corollary}[Answer to Remark 2.8]\label{cor:answer}
There is no Banach lattice \(E\) such that \(E^*\) has cotype
\(2\le p<\infty\) and \(E\) fails the disjoint property of order \(p\).
\end{corollary}

\begin{proof}
This is exactly Theorem \ref{thm:dual-cotype-disjoint-property}.
\end{proof}

\section{Verification and novelty check}

The proof has no computational dependency.  The crop was rendered from the
arXiv PDF with \texttt{code/crop\_open\_problem.py}.

For duplicate and novelty checking, the run indexes were searched for
\texttt{2309.03775}, the source title, \texttt{disjoint property of order},
\texttt{dual cotype}, and \texttt{weakly p-summable}.  No previous run packet
for this question was found.  Web searches on 2026-06-20 for the source
question, for \texttt{E* has cotype} together with \texttt{disjoint property
of order p}, and for close phrases involving bounded disjoint weakly
\(p\)-summable sequences did not locate a separate later resolution.  The
proof above is a short consequence of the source paper's Lemma 2.4 together
with the standard finite-cotype/order-continuity observation proved in
Lemma \ref{lem:cotype-oc}; human review should decide whether to describe the
result as a new observation or as an implicit consequence of already available
tools.

\section{Human review recommendation}

Recommended review focus:
\begin{itemize}
\item Check that the non-order-continuity criterion used in
  Lemma \ref{lem:cotype-oc} is the intended Banach-lattice norm criterion.
\item Check that Lemma 2.4 of \cite{botelho-garcia-miranda} is applied with
  exactly the same cotype exponent \(p\).
\item Check that the weakly null sequence plus its limit is accepted as a
  weakly compact set in the weak topology.
\end{itemize}

If these standard points are accepted, the source question is fully settled
with a negative answer.

\begin{thebibliography}{9}

\bibitem{botelho-garcia-miranda}
G. Botelho, L. A. Garcia, and V. C. C. Miranda,
\emph{Disjoint \(p\)-convergent operators and their adjoints},
arXiv:2309.03775, 2023.

\bibitem{aliprantis-burkinshaw}
C. D. Aliprantis and O. Burkinshaw,
\emph{Positive Operators},
Springer, Dordrecht, 2006.

\bibitem{meyer-nieberg}
P. Meyer-Nieberg,
\emph{Banach Lattices},
Springer-Verlag, Berlin, 1991.

\end{thebibliography}

\end{document}
